%Generated with ChatGPT from CliffTaubes's Notes
\documentclass[12pt,letterpaper]{article}
\usepackage{amsmath, amssymb, amsthm}

%%
%% Required packages:
%%
\usepackage{etex}
\usepackage{amsmath}
\usepackage{amssymb}
\usepackage{amstext}
\usepackage{amsthm}
\usepackage{fancyhdr,fancybox}
\usepackage{mathrsfs} % need for \mathscr command
\usepackage{multicol}
\usepackage[normalem]{ulem}
%\usepackage{pictex,pictexplus}
% \usepackage{pictexwd}
\usepackage{rotating}
\usepackage{url}
\usepackage{xfrac}
\usepackage{booktabs}
\usepackage{color,soul}
\usepackage{comment}

\usepackage{hyperref}
\hypersetup{
    colorlinks=true,     % false: boxed links; true: colored links
    linkcolor=blue,      % color of internal links
    citecolor=blue,      % color of links to bibliography
    filecolor=blue,      % color of file links
    urlcolor=blue        % color of external links
}


\usepackage{tikz}
\usetikzlibrary{backgrounds,calc}

%%
%% Common formatting commands
%%
\setlength{\parindent}{0in}
\setlength{\textwidth}{7in}
\setlength{\evensidemargin}{-0.25in}
\setlength{\oddsidemargin}{-0.25in}
\setlength{\parskip}{.5\baselineskip}
\setlength{\topmargin}{-0.5in}
\setlength{\textheight}{9in}

%               Problem and Part
\newcounter{problemnumber}
\newcounter{partnumber}
\newcounter{subpartnumber}
\newcommand{\Problem}{%
  \refstepcounter{problemnumber}%
  \setcounter{partnumber}{0}%
  \setcounter{subpartnumber}{0}%
  \item[\makebox{\hfil\textbf{\theproblemnumber.}\hfil}]%
  }
\newcommand{\InProblem}[2][1.6]{%
  \refstepcounter{problemnumber}%
  \setcounter{partnumber}{0}%
  \setcounter{subpartnumber}{0}%
  \textbf{\theproblemnumber.}\ \ \parbox[t]{#1in}{#2}%
}
\newcommand{\InSmallProblem}[1]{\InProblem[1.05]{#1}}
\newcommand{\NoProblem}[1]{%
  \mbox{\hfil\textbf{#1}\hfil}%
  }
\newcommand{\UnProblem}{%
  \refstepcounter{problemnumber}%
  \setcounter{partnumber}{0}%
  \setcounter{subpartnumber}{0}%
  \mbox{\hfil\textbf{\theproblemnumber.}\hfil}%
  }
\newcommand{\InFlexProblem}[1]{%
  \refstepcounter{problemnumber}%
  \setcounter{partnumber}{0}%
  \setcounter{subpartnumber}{0}%
  \textbf{\theproblemnumber.}\ \ #1%
}
%%  Part:
\newcommand{\Part}{%
  \renewcommand{\thepartnumber}{(\alph{partnumber})}%
  \refstepcounter{partnumber}
  \setcounter{subpartnumber}{0}%
  \item[(\alph{partnumber})]%
}
\newcommand{\InPart}[2][1.75]{%
  \renewcommand{\thepartnumber}{(\alph{partnumber})}%
  \refstepcounter{partnumber}%
  \setcounter{subpartnumber}{0}%
  (\alph{partnumber})\ \ \parbox[t]{#1in}{#2}%
}
\newcommand{\UnPart}{%
  \refstepcounter{partnumber}%
  \renewcommand{\thepartnumber}{(\alph{partnumber})}%
  \setcounter{subpartnumber}{0}%
  (\alph{partnumber})%
}
\newcommand{\InLargePart}[1]{\InPart[3.05]{#1}}
\newcommand{\InFullPart}[1]{\InPart[6.2]{#1}}
\newcommand{\InSmallPart}[1]{\InPart[1.05]{#1}}
%% SubPart:
\newcommand{\SubPart}{%
  \refstepcounter{subpartnumber}%
  \item[(\roman{subpartnumber})]%
  }
\newcommand{\InSubPart}[2][2.25]{%
  \stepcounter{subpartnumber}%
  (\roman{subpartnumber})\ \ \parbox[t]{#1in}{#2}%
}
\newcommand{\InSmallSubPart}[1]{\InSubPart[1.5]{#1}}
\newcommand{\InTinySubPart}[1]{%
  \stepcounter{subpartnumber}%
  (\roman{subpartnumber})\ \ \makebox{#1}%
}
\newcommand{\InMySubPart}[2][2.25]{%
  \stepcounter{subpartnumber}%
  \parbox[t]{#1in}{\makebox[0.3in][l]{(\roman{subpartnumber})}\ \ {#2}}%
}
\newcommand{\TrueFalse}{\textbf{T}\ \ \textbf{F}\ \ }
\newcommand{\TFPart}[1]{% 
  \stepcounter{partnumber}%
  \setcounter{subpartnumber}{0}%
  \item[(\alph{partnumber})] \TrueFalse\parbox[t]{5.5in}{#1}}

\newcommand{\Points}[1]{(#1 points) \ }
\newcommand{\Point}[1]{(#1 point) \ }


%% For Answers:
\newcounter{answernumber}
\newcommand{\Answer}{%
  \refstepcounter{answernumber}%
  \setcounter{partnumber}{0}%
  \setcounter{subpartnumber}{0}%
  \item[\makebox{\hfil\textbf{\theanswernumber.}\hfil}]%
  }


% Setting up the headers for the first page
%\chead{} % will default to what is typed in the file.
\fancyhead[L]{\textsc{Math 22a}}
\fancyhead[R]{\textsc{Fall 2024}}
\fancyfoot[R]{
	\vspace*{\fill}\footnotesize\textcopyright{} \emph{Harvard Math 22a}	
}
\renewcommand{\headrulewidth}{0pt}

%\fancyhead[C]{}
%	\chead{}	
%	\lhead{}
%	\rhead{}
%	\cfoot{}


% Putting copyright on the footer of each page
%\fancypagestyle{secondpage}{
%	\fancyhead[C]{TESTing again} % change this to be blank
%		%\chead{}	
%	\fancyhead[L]{bears} % change this to be blank
%	\fancyhead[R]{dogs} % change this to be blank
%	\fancyfoot[R]{ %keeps the footer the same
%		\vspace*{\fill}\footnotesize\textcopyright{} \emph{Harvard Math 22a}	
%	}
%}
%\renewcommand{\headrulewidth}{0pt}
%\pagestyle{fancy}
\pagestyle{empty}


%\lhead{}
%\rhead{}
%\rfoot{\vspace*{\fill}\footnotesize\textcopyright{} \emph{Harvard Math 22a}}
%\renewcommand{\headrulewidth}{0pt}
%\pagestyle{fancy}




%%
%% Common shortcut commands:
%%
\newcommand{\ds}{\displaystyle}
\newcommand{\sgn}{\operatorname{sgn}}
\renewcommand{\Pr}{\operatorname{P}}
\newcommand{\CPr}[3][{}]{\Pr_{\text{#1}}\left(#2\,|\,#3\right)}

\newcommand{\C}{\mathbb{C}}
\newcommand{\R}{\mathbb{R}}
\newcommand{\Z}{\mathbb{Z}}
\newcommand{\N}{\mathbb{N}}
\newcommand{\Q}{\mathbb{Q}}
\newcommand{\D}{\mathbb{D}}

\newcommand{\rref}{\operatorname{rref}}
\newcommand{\im}{\operatorname{im}}
\newcommand{\rank}{\operatorname{rank}}
\newcommand{\diag}{\operatorname{diag}}
\newcommand{\Span}{\operatorname{span}}
%\renewcommand{\vec}[1]{\mathbf{#1}}
\newcommand{\charpoly}{\mathscr{P}}
\newcommand{\nicefrac}[2]{\sfrac{#1}{#2}} % using xfrac package
\newcommand{\topology}{\mathcal{T}}
\newcommand{\Inn}{\operatorname{Inn}}

\newcommand{\blankbox}[2]{\fbox{\rule{#1}{0in}\rule{0in}{#2}}}

\newenvironment{myindentpar}[1]%
 {\begin{list}{}%
         {\setlength{\leftmargin}{#1}
          \setlength{\rightmargin}{\leftmargin}}%
         \item[]%
 }
 {\end{list}}
\newtheorem{theorem}{Theorem}
\newtheorem*{NStheorem}{Theorem}
\newtheorem{cor}[theorem]{Corollary}
\newtheorem*{claim}{Claim}
\newtheorem{defn}{Definition}

\newenvironment{amatrix}[1]{%
  \left(\begin{array}{@{}*{#1}{c}|c@{}}
}{%
  \end{array}\right)
}

%--------grstep
% For denoting a Gauss' reduction step.
% Use as: \grstep{\rho_1+\rho_3} or \grstep[2\rho_5 \\ 3\rho_6]{\rho_1+\rho_3}
\newcommand{\grstep}[2][\relax]{%
   \ensuremath{\mathrel{
       {\mathop{\longrightarrow}\limits^{#2\mathstrut}_{
                                     \begin{subarray}{l} #1 \end{subarray}}}}}}
\newcommand{\swap}{\leftrightarrow}






\chead{\large\textbf{Symmetric Matrices and the Spectral Theorem, $\vec\S$7.1}}

\begin{document}
\thispagestyle{fancy}

Recall from previously:

\begin{itemize}
\item[(a)] Definition of \textbf{inner product space}
\begin{enumerate}
    \item A function on $V \times V$: $\langle u, v \rangle$
    \item This is symmetric, linear in both $u$ and $v$.
    \item $\langle u, u \rangle > 0$ if $u \neq 0$.
\end{enumerate}

\item[(b)] \textbf{Examples}
\begin{enumerate}
    \item $\mathbb{R}^n$ and its subspaces with $\langle , \rangle = $ dot product.
    \item Polynomials of degree $n$ and evaluation at $n+1$ or more points:
    \[
    \langle P, Q \rangle = P(t_0)Q(t_0) + \cdots + P(t_n)Q(t_n).
    \]
    With regards to $\langle P, P \rangle > 0$, a polynomial that vanishes at $n+1$ points is identically zero since it has at most $n$ roots.
    \item Continuous functions on $[a, b]$ with inner product
    $
    \langle f, g \rangle = \int_a^b f(t)g(t) \, dt.
    $
\end{enumerate}

\item[(c)] \textbf{Length, distance, and orthogonality}
\begin{enumerate}
    \item $\text{dist}(u, v) = \|u - v\|.$
    \item Length of $u$ is $\|u\|.$
    \item $v$ is said to be orthogonal to $u$ if $\langle v,u \rangle = 0.$
    \item Orthogonal complements of subspaces.
    \item Orthogonal projection as before.
\end{enumerate}

\item[(d)] \textbf{Gram-Schmidt}.
This is the same as before for finite or countably infinite bases.

\item[(e)] \textbf{Best approximation}.
If $W \subset V$ is a subspace with an orthogonal basis:
\begin{enumerate}
    \item Orthogonal projection of $u$ onto a subspace $W \subset V$ with an orthogonal basis is as before:
    \[
    \text{proj}_W(u) = c_1v_1 + \cdots + c_kv_k,    
    \quad \mbox{ where} \quad c_k = \frac{\langle v_k,u \rangle}{\|v_k\|^2}.\]
    \item $\text{proj}_W(u)$ is the closest vector in $W$ to $u$ (it minimizes the function $v \mapsto \|u - v\|^2$ for $v \in W$).
\end{enumerate}

\item[(f)] \textbf{Inequalities}
\begin{enumerate}
    \item Cauchy-Schwarz inequality: \qquad
    $
    |\langle u, v \rangle| \leq \|u\| \|v\|.
    $
    \item Triangle inequality:
    \qquad $    \|u - v\| \leq \|u\| + \|v\|.
    $
    \item Angle between $u$ and $v$: 
    \qquad $
    \langle u, v \rangle = \|u\| \|v\| \cos(\theta).
    $
\end{enumerate}
\end{itemize}

\begin{defn}
An $n \times n$ matrix $A$ is said to be symmetric when $A^T = A.$
\end{defn}

\begin{enumerate}

\Problem True or false: A diagonal matrix is symmetric.
\vspace{1in}



\Problem
Prove the following theorem using the outline below.

\textbf{Transpose Theorem:}
If $u$ and $v$ are vectors in $\mathbb{R}^n$ and $A$ is an $n \times n$ matrix, then $u \cdot (Av) = (A^Tu) \cdot v.$

\begin{enumerate}
    \item[a)] Take a basis $E_{ij}$ for $n \times n$ matrices with $E_{ij}$ having all entries zero except for the $i,j$ entry which is $1.$ Verify the theorem for $E_{ij}.$\vfill
    \item[b)] Since the equation $u \cdot (Av) - (A^Tu) \cdot v = 0$ is a linear equation for the components of $A$ for any fixed choice of $u$ and $v$, it is sufficient to verify the equation for each basis element $E_{ij}.$ (Justify this statement!)\vfill
\end{enumerate}

\newpage


\Problem True or false: An orthogonally diagonalizable matrix is symmetric, where orthogonally diagonalizable is defined as follows:

\begin{defn}
An $n \times n$ matrix $A$ is said to be orthogonally diagonalizable if there is a diagonal matrix $D$ and orthogonal matrix $P$ with $P^T = P^{-1}$ such that $A = PDP^{-1}.$
\end{defn}


\vfill


\textbf{Theorem}:
A matrix is orthogonally diagonalizable if and only if it is symmetric.
\begin{enumerate}
    \item[a)] Thus, $A$ has $n$ real eigenvalues counting multiplicity.
    \item[b)] The eigenvectors form a complete, orthogonal basis.
\end{enumerate}
\vfill


\Problem Find the eigenvalues and then an orthonormal basis of eigenvectors for each of the matrices below:
\begin{enumerate}
    \item[a)] 
    $
    \begin{pmatrix}
    a & b \\
    b & c
    \end{pmatrix}
    $\vfill
    \item[b)] 
    $
    \begin{pmatrix}
    1 & 0 & 1 \\
    0 & -1 & 0 \\
    1 & 0 & 1
    \end{pmatrix}
    $\vfill
    \item[c)] 
    $
    \begin{pmatrix}
    1 & 0 & 1 \\
    0 & -1 & 0 \\
    1 & 0 & -1
    \end{pmatrix}
   $\vfill
\end{enumerate}


\newpage

\Problem True or false: If $A$ is any given $m \times n$ matrix, then $A^TA$ is a symmetric, $n \times n$ matrix.
\vfill

\Problem
True or false: If $A$ is any given $m \times n$ matrix, then the eigenvalues of $A^TA$ are the same as those of $AA^T.$
\vfill

\Problem
Let $A$ denote the $3 \times 1$ matrix
$
\begin{pmatrix}
1 \\
2 \\
3
\end{pmatrix}.
$
\begin{enumerate}
    \item[a)] Find the eigenvalues of $A^TA.$\vfill
    \item[b)] Find the eigenvalues of $AA^T.$
\end{enumerate}

\vfill

\newpage
\Problem
Let $u_1$ and $u_2$ denote vectors in $\mathbb{R}^3.$ Let $A = u_1u_1^T + u_2u_2^T.$

\begin{enumerate}
    \item[a)] True or false: The matrix $A$ always has $0$ for one of its eigenvalues.\vfill
    \item[b)] True or false: The eigenvalue $0$ has multiplicity $1$ if and only if $u_1$ and $u_2$ are not linearly independent.\vfill
\end{enumerate}

\Problem \textbf{Spectral theorem}
Let $A$ be a symmetric matrix. Fix an orthonormal basis $\{v_k\}_{k=1,\dots,n}$ of eigenvectors. Then
\[
A = \lambda_1 v_1v_1^T + \cdots + \lambda_nv_nv_n^T.
\]

Let $A$ and $\{v_1, \dots, v_n\}$ be as described in the spectral theorem. What is the matrix $A$ when written using this basis?
\vfill

\end{enumerate}

\begin{comment}
%TODO
% -Format the below and check the answers

\newpage
\chead{\Large\textbf{Symmetric matrices and spectral theorem - Answers / Solutions}}%
\lhead{}
\rhead{}
\thispagestyle{fancy}



\section*{1. True}
Since the transpose flips entries across the diagonal, any matrix $A$ has the same diagonal elements as its transpose.

\section*{2.}
\begin{enumerate}
    \item[a)] $u \cdot (E_{ij} v) = u_iv_j.$ Meanwhile $E_{ij}^T = E_{ji}$ and so $(E_{ji}u) \cdot v = v_ju_i.$
    \item[b)] Since $v \cdot (A + A^T)u = v \cdot (Au + A^Tu)$ by linearity of matrix multiplication. Meanwhile, $v \cdot (Au + A^Tu) = v \cdot Au + v \cdot A^Tu$ by linearity of the dot product with respect to either the left-hand vector or the right-hand vector.
\end{enumerate}

\section*{3. True}
If $A = UDU^{-1}$ and $U$ is orthogonal, then $A = UDUT$ and its transpose is thus $(U^T)^T D U^T$, which is $A$ since $(U^T)^T = U$ and $D^T = D$ and $U^T = U^{-1}.$

\section*{4.}
\begin{enumerate}
    \item[a)] Eigenvalues are $\lambda = \frac{1}{2} (a + c) \pm \frac{1}{2} \sqrt{(a - c)^2 + 4b^2}.$ Whichever version of $\lambda$, a corresponding eigenvector is $v = \begin{pmatrix} b \\ \lambda - a \end{pmatrix}.$
    \item[b)] Eigenvalues are $-1, 0, 2$ since the characteristic polynomial is $\lambda(1 + \lambda)(2 - \lambda).$ The corresponding eigenvectors are:
    \[
    \begin{pmatrix}
    0 \\
    1 \\
    0
    \end{pmatrix},
    \begin{pmatrix}
    1 \\
    0 \\
    -1
    \end{pmatrix},
    \begin{pmatrix}
    1 \\
    0 \\
    1
    \end{pmatrix}.
    \]
    \item[c)] Eigenvalues are $-1, -\sqrt{2},$ and $\sqrt{2}$ since the characteristic polynomial is $(1 + \lambda)(2 - \lambda^2).$ The corresponding eigenvectors are:
    \[
    \begin{pmatrix}
    0 \\
    1 \\
    0
    \end{pmatrix},
    \begin{pmatrix}
    1 \\
    0 \\
    -1 - \sqrt{2}
    \end{pmatrix},
    \begin{pmatrix}
    1 \\
    0 \\
    -1 + \sqrt{2}
    \end{pmatrix}.
    \]
\end{enumerate}

\section*{5. True}
Because $(A^TA)^T = A^T(A^T)^T = A^TA.$

\section*{6. False}
See the answer to \#7 on this worksheet. But it is true if $A$ is a square matrix. For example, similar matrices have the same eigenvalues; and if $A$ is invertible, then $A^TA$ is similar to $AA^T$ because $A^{-1}(AA^T)A = A^TA.$

\section*{7.}
\begin{enumerate}
    \item[a)] $A^TA$ is the $1 \times 1$ matrix with entry $14.$
    \item[b)] $AA^T = \begin{pmatrix} 1 & 2 & 3 \\ 2 & 4 & 6 \\ 3 & 6 & 9 \end{pmatrix}$ which has eigenvalues $0$ (with multiplicity 2) and $14.$
\end{enumerate}

\section*{8.}
\begin{enumerate}
    \item[a)] True, and an eigenvector with eigenvalue $0$ is any non-zero vector orthogonal to the span of $u_1$ and $u_2.$
    \item[b)] If $u_1$ and $u_2$ are linearly independent, then the orthogonal complement to their span has dimension 1. If they are not linearly independent, then it has dimension 2.
\end{enumerate}

\section*{9.}
The matrix $A$ is the diagonal matrix with the eigenvalues $\lambda_1, \dots, \lambda_n$ on the diagonal.

\end{comment}

\end{document}


